\documentclass{article}

\usepackage[utf8]{inputenc}
\usepackage[turkish, shorthands = :!]{babel}
\usepackage[T1]{fontenc}
\usepackage{lmodern}

\usepackage{fancyhdr, lipsum, mathtools}

\usepackage{siunitx}
\sisetup{output-decimal-marker = {,}}
\sisetup{group-separator = {.}}
\sisetup{group-minimum-digits = 4}

\renewcommand{\headrulewidth}{0 pt}
\pagestyle{fancy}
\fancyhf{}
\fancyhead[C]{\thepage}

\begin{document}

%%%%%%% SAYFA 1 %%%%%%%

\lipsum[1][1-3]
% sergilenen (displayed) bir denklemden önce kesinlikle boş satır bırakmamalı (illa bırakılacaksa o satıra en azından bir yüzde işareti konmalı) ki denklem öncesi ve sonrası boşluklar dengeli olsun
\[ i = \sqrt{-1} \] % çeşitli sebeplerden ötürü \[...\] yerine $$...$$ kullanmak önerilmiyor
% bu satırı boş bırakmak dikey boşluğu değiştirmiyor, sadece müteakip cümlenin paragraf başı olup olmamasını belirliyor
\lipsum[1][4-6]

\bigskip
\hrule
\bigskip

\lipsum[1][1-3]

\[ i = \sqrt{-1} \]

\lipsum[1][4-6]

\bigskip
\hrule
\bigskip

\lipsum[1][1-3]
\begin{equation} \label{eq:imaginary}
i = \sqrt{-1}
\end{equation}
\lipsum[1][4-6]

Denklem~\eqref{eq:imaginary}...

\bigskip
\hrule
\bigskip

% yanlış
\lipsum[1][1-3]
\[ i = \sqrt{-1} \]
\[ ax^2 + bx + c = 0 \]
\lipsum[1][4-6]

\bigskip
\hrule
\bigskip

% doğru
\lipsum[1][1-3]
\begin{gather*}
i = \sqrt{-1} \\ 
ax^2 + bx + c = 0 % buraya \\ işareti konmamalı
\end{gather*}
\lipsum[1][4-6]

\newpage

%%%%%%% SAYFA 2 %%%%%%%

\lipsum[1][1-3]
\begin{multline*}
\Big( 1 - \frac{w_i}{W} \Big) \Big( 1 - \frac{w_i}{W - w_N} \Big) \Big( 1 - \frac{w_i}{W - w_N - w_{N - 1}} \Big) \\
\leq q_i \leq \Big( 1 - \frac{w_i}{W} \Big) \Big( 1 - \frac{w_i}{W - w_1} \Big) \Big( 1 - \frac{w_i}{W - w_1 - w_2} \Big).
\end{multline*}
\lipsum[1][4-6]

\bigskip
\hrule
\bigskip

\lipsum[1][1-3]
\begin{align*}
P(X < \num{2,99} & \text{ veya } X > \num{3,01}) = P(X < \num{2,99}) + P(X > \num{3,01}) \\
&= P\left(\frac{X-\mu}{\sigma} < \frac{\num{2,99}-\mu}{\sigma}\right) + P\left(\frac{X-\mu}{\sigma} > \frac{\num{3,01}-\mu}{\sigma}\right) \\
&= P\left(Z < \frac{\num{2,99}-\mu}{\sigma}\right) + P\left(Z > \frac{\num{3,01}-\mu}{\sigma}\right) \\
&= P(Z < -2) + P(Z > 2) \\
&= \num{0,0456}.
\end{align*}
\lipsum[1][4-6]

\newpage

%%%%%%% SAYFA 3 %%%%%%%

\lipsum[1][1-3]
\begin{align*} 
y	& = (x - a)(x - b) \\
	& = x^2 - (a + b)x + ab
\end{align*}
\lipsum[1][4-6]

\bigskip
\hrule
\bigskip

\lipsum[1][1-3]
\begin{align*}
4x - y + 2z		& = -1 \\
-x + (\ln 2)y - 3z	& = 0 % buraya \\ işareti konmamalı
\end{align*}
\lipsum[1][4-6]

\bigskip
\hrule
\bigskip

\lipsum[1][1-3]
\begin{align*}
x		& = y		&	w		& = z \\
2x		& = -y		&	3w		& = \frac{1}{2}z \\
-4 + 5x	& = 2 + y	&	w + 2	& = -1 + w       
\end{align*}
\lipsum[1][4-6]

\newpage

%%%%%%% SAYFA 4 %%%%%%%

\lipsum[1][1-3]
\begingroup
\sisetup{group-separator = {}}
\setlength{\jot}{0.75 em} % alt alta kesirli denklemlerde satırlar arası boşluğu artırmak için (varsayılan değeri 3 pt)
\begin{align*}
P(B_2|A)
	&= \frac{P(B_2)P(A|B_2)}{\sum_{i=1}^{3} P(B_i)P(A|B_i)} \\
	&= \frac{\num{0,45} \cdot \num{0,03}}{\num{0,3} \cdot \num{0,02} + \num{0,45} \cdot \num{0,03} + \num{0,25} \cdot \num{0,02}} \\
	&= \frac{\num{0,0135}}{\num{0,0245}} \\
	&= \frac{27}{49}.
\end{align*}
\endgroup
\lipsum[1][4-6]

\end{document}